\documentclass[../sparc.tex]{subfiles}
\graphicspath{{\subfix{../images/}}}
\begin{document}

\newpage
%%%%%%%%%%%%%%%%%%%%%%%%%%%%%%%%%%%%%%%%%%%%%%%%%%%%%%%%%%%%%%%%%%%%%%%%%%%%%%%%
\section{Creating a game engine}
\index{Game development!Game engine}

%%%%%%%%%%%%%%%%%%%%%%%%%%%%%%%%%%%%%%%%%%%%%%%%%%%%%%%%%%%%%%%%%%%%%%%%%%%%%%%%
\subsection{What is a ``Game engine''?}

During game development programmers, accumulate a set of reusable procedures and
data types (structures) that were created specifically to make the working on
the game easier.  Using those specialized tools we can quickly develop the game
logic without the need to deal with the low-level mechanisms, such as drawing
pixels on a screen, moving objects around a map etc.

When a team of developers faces with the need to create a new version of a game
or even with development of a new game, it usually turns out that basic problems
have already been solved in the previous game.  It's not surprising as working
with a game map, playing sounds and calculating physics exist in many games.
Why repeating the work that is already done, when we can use the existing
procedures and data structures?

Thus we come to the idea of making our own \emph{game engine} -- a library of
procedures and types created specifically for game development.  Even if we use
our engine for only one game it will give us an advantage in terms of and
ordering of in modularity and code organization.

At the most basic level development of a game engine boils down to moving our
data types and procedures into separate files.  This is the next level of
modularity: first we separated the logic into procedures; now we will group
those procedures into separate files according to their purpose.

To figure out how to split our project into files we must understand what the
process of building our project into a single executable file entails.

%%%%%%%%%%%%%%%%%%%%%%%%%%%%%%%%%%%%%%%%%%%%%%%%%%%%%%%%%%%%%%%%%%%%%%%%%%%%%%%%
\subsection{Stages of application building}
\index{Programming!Application build process!Compiler}

\newglossaryentry{compiler}{name=compiler, description={A program that
    translates a text written in one programming language (source code) into
    other (target code.)}}

Up to this point we discussed application build process as a single, indivisible
process: we click a button Arduino IDE interface and then computer immediately
starts performing some operations on our project's code, and then it uploads the
program into an Arduino and finally we get long-awaited blinking LED.  But not
everything is as easy as it seems at the first glance.

Of course we have already saw already encountered situations when a project
refused to build due to syntax errors in our code: forgotten parenthesis and
semicolons make themselves felt at the early stages of the build process.  In
this case we could see error messages (that sometimes freaking out novice
programmers) in Arduino IDE.  Now it is time to lift the curtain of secrecy from
compilation process and look inside.

In the beginning computer starts its work with the source code that we wrote.
It's important to understand that a computer (or, more precisely, its processor)
does not understand our programming language and a special ``translator'' from
the intermediate language (that is, programming language) to the low-level
machine commands is required.  The \gls{compiler} acts as the ``translator''
that translates text written in one programming language (source code) into
other (target code.)  The target code can be either a \emph{machine instruction}
language (a processor understood language) or a non-machine language.  In our
case the target language is the machine instruction language.

But compilation is a quite complicated process that can be divided into several
main steps as is shown on the fig. \ref{fig:build-process}.

\figureBuildSteps{en}

Let's discuss those steps in order.

\subsubsection{Preprocessor}
\index{Programming!Application build process!Preprocessor}

First of all the source code (\texttt{main.ino} file) is being processed by the
\emph{preprocessor} -- a special program that is a part of the compiler.  On this
stage commentaries are being deleted and special \emph{preprocessor directives}
(sometimes called \emph{preprocessor commands}) are being handled.  In C/C++
language the preprocessor language is a separate programming language with quite
primitive set of commands and describes operations that preprocessor must do
with the source code.

Essentially preprocessor operates with tree basic tools: code deletion, copying
and inserting.

For example, as we discussed above, commentaries (all the lines that start with
two right slashes, or lines that are in between \mintinline{cpp}{/*...*/}) are
being deleted and do not proceed to the subsequent stages of building.

Other popular preprocessor directives that we can encounter are
\mintinline{cpp}{#include} and \mintinline{cpp}{#define} commands.

\mintinline{cpp}{#define} command is usually used in the following way:

\begin{listing}[H]
  \begin{minted}[highlightlines={1, 8, 10}]{cpp}
    #define DELAY_TIME 100
    void setup() {
      pinMode(2, OUTPUT);
    }

    void loop() {
      digitalWrite(2, HIGH);
      delay(DELAY_TIME);
      digitalWrite(2, LOW);
      delay(DELAY_TIME);
    }
  \end{minted}
  \label{listing:preprocessor-define-example-1}
  \caption{An example of using a preprocessor directive for specifying a
    constant.}
\end{listing}

In the listing \ref{listing:preprocessor-define-example-1} we can see an example
of \mintinline{cpp}{#define} use for specifying a named constant
\mintinline{cpp}{DELAY_TIME} that sets the time of delay between the turning on
and off the LED.  In the \mintinline{cpp}{loop} procedure this constant is being
substituted as the value of delay in the \mintinline{cpp}{delay} procedure call.

When the preprocessor encountering a constant defined with
\mintinline{cpp}{#define} before, it replaces its name with the factual value of
the constant.  It's important to understand here that preprocessor performs
almost no checks regarding the correctness of the result of the constant
substitution.  Thus in the case of a mistakenly specified constant, the
preprocessor will substitute its value all over the code where the constant is
mentioned, but on the later stages of the compilation process an error will
occur.

In addition we can see on the first line of the listing
\ref{listing:preprocessor-define-example-1} that we don't use already familiar
to us data types, and we even don't put an ubiquitous semicolon at the end;
moreover, terminating the a \mintinline{cpp}{#define} line with a semicolon will
subsequently lead to an error.

The question arises: should we use \mintinline{cpp}{#define} when we have
\mintinline{cpp}{const} keyword?  We might as well do the following:

\begin{listing}[H]
  \begin{minted}[highlightlines={1}]{cpp}
    const int DELAY_TIME = 100;
    void setup() {
      pinMode(2, OUTPUT);
    }

    void loop() {
      digitalWrite(2, HIGH);
      delay(DELAY_TIME);
      digitalWrite(2, LOW);
      delay(DELAY_TIME);
    }
  \end{minted}
  \label{listing:preprocessor-define-example-2}
  \caption{An example of specifying a regular constant.}
\end{listing}

Indeed, the approach that is shown in the listing
\ref{listing:preprocessor-define-example-2} has a lot of advantages comparing to
the listing \ref{listing:preprocessor-define-example-1} where we used
\mintinline{cpp}{#define}.  For example, the constant that we defined with
\mintinline{cpp}{const} is typed (uses \mintinline{cpp}{int} data type) and all
necessary syntax correctness checks are performed when compiling such code
(unlike \mintinline{cpp}{#define}.)  For this reason in most of the cases usage
of regular constants are preferable to usage of \mintinline{cpp}{#define}.

However, we shouldn't rule out using preprocessor directives as there are cases
where using preprocessor is acceptable and even required.  The thing is that the
preprocessor performs operations specifically with on program's source code and
thus the later compiling stages receive already processed code.

\end{document}
